\subsection*{الف}
برای بررسی \lr{orthonormal} بودن، ضرب داخلی‌ها را حساب می‌کنیم:
\begin{flalign*}
s_0^T s_0 &= \left(\frac{1}{\sqrt{2}}\begin{bmatrix}1\\[2pt]1\end{bmatrix}\right)^T
           \left(\frac{1}{\sqrt{2}}\begin{bmatrix}1\\[2pt]1\end{bmatrix}\right)
= \frac{1}{2}(1^2+1^2)=1, \\
s_1^T s_1 &= \left(\frac{1}{\sqrt{2}}\begin{bmatrix}1\\[2pt]-1\end{bmatrix}\right)^T
           \left(\frac{1}{\sqrt{2}}\begin{bmatrix}1\\[2pt]-1\end{bmatrix}\right)
= \frac{1}{2}(1^2+(-1)^2)=1, \\
s_0^T s_1 &= \left(\frac{1}{\sqrt{2}}\begin{bmatrix}1\\[2pt]1\end{bmatrix}\right)^T
           \left(\frac{1}{\sqrt{2}}\begin{bmatrix}1\\[2pt]-1\end{bmatrix}\right)
= \frac{1}{2}(1\cdot 1 + 1\cdot (-1))=0. &&
\end{flalign*}
پس هر دو بردار نُرمِ ۱ دارند و نسبت به هم عمود هستند؛ بنابراین پایه‌ها \lr{orthonormal} هستند.

\subsection*{ب}
ماتریس تبدیل را از کنار هم گذاشتن پایه‌ها می‌سازیم (به عنوان ستون‌ها):
\begin{flalign*}
A = [s_0\ s_1] = \frac{1}{\sqrt{2}}
\begin{bmatrix}
1 & 1\\
1 & -1
\end{bmatrix}, \qquad A^T A = I \;\; \Rightarrow \;\; A^{-1}=A^T. &&
\end{flalign*}
تبدیلِ ماتریسِ $F$ در این پایه برابر است با:
\[
T = A^T F A.
\]
محاسبه:
\begin{flalign*}
F A
&=
\begin{bmatrix}
20 & 63\\
21 & 128
\end{bmatrix}
\cdot
\frac{1}{\sqrt{2}}
\begin{bmatrix}
1 & 1\\
1 & -1
\end{bmatrix}
=
\frac{1}{\sqrt{2}}
\begin{bmatrix}
83 & -43\\
149 & -107
\end{bmatrix},\\[6pt]
T
&=A^T (F A)
=
\frac{1}{\sqrt{2}}
\begin{bmatrix}
1 & 1\\
1 & -1
\end{bmatrix}
\cdot
\frac{1}{\sqrt{2}}
\begin{bmatrix}
83 & -43\\
149 & -107
\end{bmatrix}
=
\frac{1}{2}
\begin{bmatrix}
232 & -150\\
-66 & 64
\end{bmatrix}
=
\begin{bmatrix}
116 & -75\\
-33 & 32
\end{bmatrix}. &&
\end{flalign*}

\subsection*{ج}
چون $A$ \lr{orthogonal} است، تبدیلِ معکوس به صورت زیر است:
\[
F = A T A^T.
\]
برای اثبات بازیابیِ دقیق، مستقیماً ضرب می‌کنیم:
\begin{flalign*}
T A
&=
\begin{bmatrix}
116 & -75\\
-33 & 32
\end{bmatrix}
\cdot
\frac{1}{\sqrt{2}}
\begin{bmatrix}
1 & 1\\
1 & -1
\end{bmatrix}
=
\frac{1}{\sqrt{2}}
\begin{bmatrix}
41 & 191\\
-1 & -65
\end{bmatrix},\\[6pt]
A(TA)
&=
\frac{1}{\sqrt{2}}
\begin{bmatrix}
1 & 1\\
1 & -1
\end{bmatrix}
\cdot
\frac{1}{\sqrt{2}}
\begin{bmatrix}
41 & 191\\
-1 & -65
\end{bmatrix}
=
\frac{1}{2}
\begin{bmatrix}
40 & 126\\
42 & 256
\end{bmatrix}
=
\begin{bmatrix}
20 & 63\\
21 & 128
\end{bmatrix}
=F. &&
\end{flalign*}
پس تبدیلِ معکوس، دقیقاً ماتریس اولیه $F$ را بازسازی می‌کند.